\chapter{The parameter budget, counted by hand}\label{ch:params}

\begin{onething}
With three counting rules (a convolution has $c_{\text{in}} c_{\text{out}} k^2 + c_{\text{out}}$ numbers,
a depthwise one $c k^2 + c$, a LayerNorm $2c$) you can derive the size of Indradhanu layer by layer,
5,506,629 trainable numbers for the 165-band model, and land on the real model's figure to the last
digit; then the same rules tell you where the parameters live (blocks and decoder), where the compute
lives (the decoder), and which parameters a masked pass never trains.
\end{onething}

Every chapter of this part ended with a shape. This one ends with a count. Counting parameters is the
strongest test of understanding an architecture: if any layer in your head is the wrong size, the total
will not match, and you will know. The totals are checked against two independent sources: the
\texttt{torchinfo} table in \file{research/model_break_down/05_hyperspectral_segformer_mae.md}, and a
build of the real classes on CPU by \file{research/textbook/tools/compute_params.py}, whose output is
saved beside it as \file{compute_params_output.txt}.

\section{What a parameter is, and what it is not}\label{sec:params:what}

\begin{definition}[Parameter, buffer]
A \index{parameter}\emph{parameter} is one number that training adjusts: a single weight or bias. A
\index{buffer}\emph{buffer} is a number stored in the model file that training does not adjust by
gradient: the normalisation means and standard deviations, and BatchNorm's running statistics.
\end{definition}

Each parameter is a 32-bit float, 4 bytes, so the parameter count fixes the file size:
\[
5{,}506{,}629 \times 4 = 22{,}026{,}516 \text{ bytes} \approx 22.03 \text{ MB},
\]
which is the ``Params size (MB): 22.03'' that \texttt{torchinfo} reports.
The figure is measured: \tagM{research/model_break_down/05_hyperspectral_segformer_mae.md}.

Why bother counting? Three reasons. A count that matches proves you know every layer's shape. A count
tells you where capacity is spent, which matters when you decide what to widen or shrink
(\cref{ch:heavy}). And a count is \emph{not} a measure of cost, which is a lesson worth learning
with real numbers (\cref{sec:params:compute}).

\section{Three counting rules}\label{sec:params:rules}

\paragraph{A convolution.}\index{convolution!parameter count} \code{Conv2d(c_in, c_out, k)} has, for each of its $c_{\text{out}}$ output
channels, one kernel of $c_{\text{in}} \times k \times k$ weights (it reads every input channel over a
$k \times k$ window) and one bias:
\begin{equation}\label{eq:params:conv}
P_{\text{conv}} = c_{\text{out}}\,(c_{\text{in}} k^2 + 1) = c_{\text{in}} c_{\text{out}} k^2 + c_{\text{out}} .
\end{equation}
\paragraph{A depthwise convolution.}\index{depthwise convolution!parameter count} With \code{groups = c}, each channel has its own $k \times k$ kernel
and reads only itself, so the factor $c_{\text{in}}$ drops to 1:
\begin{equation}\label{eq:params:dw}
P_{\text{dw}} = c\,(k^2 + 1) = c k^2 + c .
\end{equation}
\paragraph{A linear layer and a LayerNorm.}\index{linear layer!parameter count}\index{LayerNorm!parameter count}\index{BatchNorm!buffers}\code{Linear(a, b)} is a $b \times a$ weight matrix plus $b$
biases; it is a convolution with $k = 1$. A LayerNorm over $c$ features learns one scale and one shift
per feature (\cref{ch:block}):
\begin{equation}\label{eq:params:lin}
P_{\text{lin}} = a b + b, \qquad P_{\text{LN}} = 2c .
\end{equation}
BatchNorm also learns $2c$ (scale and shift) and in addition keeps $2c + 1$ buffers: a running mean and a
running variance per channel and one counter of batches seen.

\begin{worked}[label={we:params:tiny}]{The rules on a tiny convolution}
\code{Conv2d(2, 3, k=2)}. Each of the 3 output channels has a kernel of $2 \times 2 \times 2 = 8$
weights and one bias. Count on paper:
\begin{align*}
\text{weights} &= 2 \times 3 \times 2^2 = 2 \times 3 \times 4 = 24,\\
\text{biases} &= 3,\\
P &= 24 + 3 = 27 .
\end{align*}
The depthwise version with 3 channels, \code{Conv2d(3, 3, k=2, groups=3)}: $3 \times 4 + 3 = 12 + 3 =
15$. The number of heads in an attention layer, the patch size and the batch size never appear in any
rule: they change shapes, not parameters.
\end{worked}

\section{The configuration being counted}\label{sec:params:config}

\begin{table}[htbp]\centering
\caption{The configuration of the 165-band Indradhanu (v0.2.0).}\label{tab:params:config}
\begin{tabular}{@{}ll@{}}\toprule
Input bands & 165 \\
Compressed width $D$ & 32 \\
Stage widths $C_1..C_4$ & 32, 64, 160, 256 \\
Reduction ratios $R_1..R_4$ & 8, 4, 2, 1 \\
Blocks per stage & 2, 2, 2, 2 \\
Mix-FFN expansion $E$ & 4 \\
Decoder width & 256 \\
Heads & 2, 2, 5, 8 (no effect on the count) \\\bottomrule
\end{tabular}
\end{table}

\section{The front end: the spectral compressor}\label{sec:params:compressor}

\Cref{ch:compressor}: a \code{Conv2d(165, 32, k=1)} and a \code{BatchNorm2d(32)}.
\begin{align*}
\text{conv} &= 165 \times 32 \times 1 \times 1 + 32 = 5{,}280 + 32 = 5{,}312,\\
\text{BatchNorm} &= 2 \times 32 = 64,\\
\text{compressor} &= 5{,}312 + 64 = 5{,}376 .
\end{align*}

\section{The encoder}\label{sec:params:encoder}

\subsection{The four patch embeddings}

Each is a convolution followed by a LayerNorm over its output width (\cref{ch:patchembed}). Stage 1
reads $D = 32$ channels with $k = 4$; stages 2--4 read the previous width with $k = 3$
(\cref{ch:encoder}).
\begin{align*}
\text{stage 1:}\quad & 32 \times 32 \times 16 + 32 = 16{,}384 + 32 = 16{,}416; & +\ 2 \times 32 &= 16{,}480\\
\text{stage 2:}\quad & 32 \times 64 \times 9 + 64 = 18{,}432 + 64 = 18{,}496; & +\ 2 \times 64 &= 18{,}624\\
\text{stage 3:}\quad & 64 \times 160 \times 9 + 160 = 92{,}160 + 160 = 92{,}320; & +\ 2 \times 160 &= 92{,}640\\
\text{stage 4:}\quad & 160 \times 256 \times 9 + 256 = 368{,}640 + 256 = 368{,}896; & +\ 2 \times 256 &= 369{,}408
\end{align*}
\[
16{,}480 + 18{,}624 + 92{,}640 + 369{,}408 = 497{,}152 .
\]

\subsection{One block, as a formula}

A block (\cref{ch:block}) of width $C$ with reduction ratio $R$ and expansion 4 contains:
\begin{itemize}
\item \textbf{ESA} (\cref{ch:esa}): four \code{Linear(C, C)} (query, key, value, output projection),
and, only when $R > 1$, a reduction \code{Conv2d(C, C, k=R, s=R)} and a \code{LayerNorm(C)};
\item \textbf{Mix-FFN}: \code{Linear(C, 4C)}, a depthwise \code{Conv2d(4C, 4C, k=3)}, \code{Linear(4C, C)};
\item \textbf{two LayerNorms} (pre-norm before each sub-layer).
\end{itemize}
Count each piece with the rules, in symbols:
\begin{align*}
\text{q, k, v, out} &= 4\,(C^2 + C) = 4C^2 + 4C,\\
\text{reduction conv} + \text{norm} &= (C \cdot C \cdot R^2 + C) + 2C = R^2 C^2 + 3C \quad (\text{if } R > 1),\\
\text{fc1} &= C \cdot 4C + 4C = 4C^2 + 4C,\\
\text{depthwise} &= 4C \cdot 9 + 4C = 40C,\\
\text{fc2} &= 4C \cdot C + C = 4C^2 + C,\\
\text{two LayerNorms} &= 2 \times 2C = 4C .
\end{align*}
Add them, collecting the $C^2$ terms ($4 + 4 + 4 = 12$) and the $C$ terms ($4 + 4 + 40 + 1 + 4 = 53$):
\begin{equation}\label{eq:params:block}
P_{\text{block}}(C, R) = 12 C^2 + 53 C + [R > 1]\,\big(R^2 C^2 + 3C\big),
\end{equation}
where $[R > 1]$ is 1 when $R > 1$ and 0 otherwise. A closed form for the whole block\index{transformer block!parameter count}.

\begin{worked}{The four stages from \texorpdfstring{\eqref{eq:params:block}}{the block formula}}\label{we:params:blocks}
Do each on paper and compare with the column-by-column count of the course in
\cref{tab:params:blockdetail}.

Stage 1, $C = 32$, $R = 8$:
\begin{align*}
12 \times 32^2 &= 12 \times 1{,}024 = 12{,}288, & 53 \times 32 &= 1{,}696,\\
8^2 \times 32^2 &= 64 \times 1{,}024 = 65{,}536, & 3 \times 32 &= 96,\\
P &= 12{,}288 + 1{,}696 + 65{,}536 + 96 = 79{,}616; & \times 2 \text{ blocks} &= 159{,}232 .
\end{align*}
Stage 2, $C = 64$, $R = 4$:
\begin{align*}
12 \times 4{,}096 &= 49{,}152, & 53 \times 64 &= 3{,}392,\\
16 \times 4{,}096 &= 65{,}536, & 3 \times 64 &= 192,\\
P &= 49{,}152 + 3{,}392 + 65{,}536 + 192 = 118{,}272; & \times 2 &= 236{,}544 .
\end{align*}
Stage 3, $C = 160$, $R = 2$:
\begin{align*}
12 \times 25{,}600 &= 307{,}200, & 53 \times 160 &= 8{,}480,\\
4 \times 25{,}600 &= 102{,}400, & 3 \times 160 &= 480,\\
P &= 307{,}200 + 8{,}480 + 102{,}400 + 480 = 418{,}560; & \times 2 &= 837{,}120 .
\end{align*}
Stage 4, $C = 256$, $R = 1$ (no reduction):
\begin{align*}
12 \times 65{,}536 &= 786{,}432, & 53 \times 256 &= 13{,}568,\\
P &= 786{,}432 + 13{,}568 = 800{,}000; & \times 2 &= 1{,}600{,}000 .
\end{align*}
The round 800,000 is a coincidence of these widths. Blocks in total:
\[
159{,}232 + 236{,}544 + 837{,}120 + 1{,}600{,}000 = 2{,}832{,}896 .
\]
Every stage figure matches \texttt{torchinfo} and the CPU build.
\end{worked}

\begin{table}[htbp]\centering
\caption{One block per stage, piece by piece (two blocks per stage in the model). These per-module
figures are exactly what the CPU build reports for stage 1 (q, k, v and proj 1,056 each, reduction
65,568, reduction norm 64, fc1 4,224, dwconv 1,280, fc2 4,128).}\label{tab:params:blockdetail}
\begin{tabular}{@{}lrrrr@{}}\toprule
Piece & Stage 1 & Stage 2 & Stage 3 & Stage 4 \\\midrule
q, k, v, out & 4,224 & 16,640 & 103,040 & 263,168 \\
reduction conv & 65,568 & 65,600 & 102,560 & -- \\
reduction norm & 64 & 128 & 320 & -- \\
fc1 & 4,224 & 16,640 & 103,040 & 263,168 \\
depthwise conv & 1,280 & 2,560 & 6,400 & 10,240 \\
fc2 & 4,128 & 16,448 & 102,560 & 262,400 \\
two LayerNorms & 128 & 256 & 640 & 1,024 \\\midrule
one block & 79,616 & 118,272 & 418,560 & 800,000 \\\bottomrule
\end{tabular}
\end{table}

Two things stand out in \eqref{eq:params:block}. The parameters grow with $C^2$: doubling the width
quadruples the block. And the reduction convolution $R^2 C^2$ is as large as all the rest of the block
at stage 1 ($65{,}536$ against $14{,}080$): a stride-8 kernel over 32 channels has $8 \times 8 \times 32
\times 32$ weights. Stages 1 and 2 have the same reduction weight count, $64 \times 32^2 = 16 \times
64^2 = 65{,}536$, because $R C$ is 256 in both.

\subsection{The stage-output norms, and the encoder total}

One LayerNorm per stage: $2 \times (32 + 64 + 160 + 256) = 2 \times 512 = 1{,}024$. The encoder:
\[
497{,}152 + 2{,}832{,}896 + 1{,}024 = 3{,}331{,}072 .
\]

\section{The decoder and the back end}\label{sec:params:decoder}

From \cref{ch:decoder}, with the rule \eqref{eq:params:conv}:
\begin{align*}
\text{projections:}\quad & 32 \times 256 + 256 = 8{,}448\\
 & 64 \times 256 + 256 = 16{,}640\\
 & 160 \times 256 + 256 = 41{,}216\\
 & 256 \times 256 + 256 = 65{,}792\\
\text{fuse:}\quad & 1{,}024 \times 256 + 256 = 262{,}144 + 256 = 262{,}400\\
\text{refine 1:}\quad & 256 \times 256 \times 9 + 256 = 589{,}824 + 256 = 590{,}080\\
\text{refine 2:}\quad & 256 \times 512 \times 9 + 512 = 1{,}179{,}648 + 512 = 1{,}180{,}160
\end{align*}
\[
8{,}448 + 16{,}640 + 41{,}216 + 65{,}792 + 262{,}400 + 590{,}080 + 1{,}180{,}160 = 2{,}164{,}736 .
\]
Pixel shuffle, bilinear interpolation and GELU have no parameters.
\code{Conv2d(32, 165, k=1)} is the decompressor:
\[
32 \times 165 + 165 = 5{,}280 + 165 = 5{,}445 .
\]

\section{The grand total, and the buffers}\label{sec:params:total}

\begin{table}[htbp]\centering
\caption{Indradhanu (165 bands, $D = 32$), every parameter. Measured by
\texttt{torchinfo}\tagM{research/model_break_down/05_hyperspectral_segformer_mae.md} and by the CPU
build\tagM{research/textbook/tools/compute_params_output.txt}; both report 5,506,629.}\label{tab:params:total}
\begin{tabular}{@{}lrr@{}}\toprule
Component & Parameters & Share \\\midrule
Spectral compressor & 5,376 & 0.10\% \\
Patch embeddings & 497,152 & 9.03\% \\
Transformer blocks & 2,832,896 & 51.44\% \\
Stage-output norms & 1,024 & 0.02\% \\
Decoder & 2,164,736 & 39.31\% \\
Spectral decompressor & 5,445 & 0.10\% \\\midrule
\textbf{Total} & \textbf{5,506,629} & 100\% \\\bottomrule
\end{tabular}
\end{table}

Add the column on paper:
\begin{align*}
5{,}376 + 497{,}152 &= 502{,}528\\
502{,}528 + 2{,}832{,}896 &= 3{,}335{,}424\\
3{,}335{,}424 + 1{,}024 &= 3{,}336{,}448\\
3{,}336{,}448 + 2{,}164{,}736 &= 5{,}501{,}184\\
5{,}501{,}184 + 5{,}445 &= 5{,}506{,}629 .
\end{align*}
\texttt{torchinfo} reports ``Total params: 5,506,629''. Exact to the last digit. (Shares: each component
divided by 5,506,629; for the blocks $2{,}832{,}896 / 5{,}506{,}629 \approx 0.5144$.)

\begin{worked}[label={we:params:buffers}]{Where the manifest's 5,507,354 comes from}
The checkpoint manifest says \code{"params": 5507354}, as the model breakdown
records\tagM{research/model_break_down/05_hyperspectral_segformer_mae.md}. The difference is
\[
5{,}507{,}354 - 5{,}506{,}629 = 725 .
\]
Count the buffers. The normaliser and denormaliser each store a mean and a standard deviation per
band:
\[
4 \times 165 = 660 .
\]
The compressor's BatchNorm keeps a running mean and running variance for each of its 32 channels and
one batch counter:
\[
32 + 32 + 1 = 65 .
\]
Together $660 + 65 = 725$, and $5{,}506{,}629 + 725 = 5{,}507{,}354$. The CPU build reports 725 buffer
numbers. Those 660 normalisation values live inside the checkpoint, which is why you cannot swap
normalisation statistics without rebuilding the model (\cref{ch:bandpipeline}).
\end{worked}

\section{Where the parameters live}\label{sec:params:where}

\begin{figure}[htbp]\centering
\begin{tikzpicture}
\begin{axis}[
  xbar, width=0.82\linewidth, height=9.2cm,
  xmin=0, xmax=1900000,
  bar width=7pt, enlarge y limits=0.04,
  xlabel={parameters},
  scaled x ticks=false,
  xticklabel style={/pgf/number format/fixed,/pgf/number format/1000 sep={,}},
  xtick={0,500000,1000000,1500000},
  symbolic y coords={decompressor,refine 2,refine 1,fuse,projections,stage norms,blocks S4,blocks S3,blocks S2,blocks S1,patch embeds,compressor},
  ytick=data, y tick label style={font=\small},
  nodes near coords, nodes near coords align={horizontal},
  every node near coord/.append style={font=\scriptsize,/pgf/number format/fixed,/pgf/number format/1000 sep={,}},
  axis x line*=bottom, axis y line*=left]
\addplot[fill=spec, draw=spec] coordinates {(5376,compressor) (5445,decompressor)};
\addplot[fill=spat, draw=spat] coordinates {(497152,patch embeds) (159232,blocks S1) (236544,blocks S2) (837120,blocks S3) (1600000,blocks S4) (1024,stage norms)};
\addplot[fill=attn, draw=attn] coordinates {(132096,projections) (262400,fuse) (590080,refine 1) (1180160,refine 2)};
\end{axis}
\end{tikzpicture}
\caption{Where Indradhanu's 5,506,629 parameters live, drawn to scale from the counts of this chapter.
Spectral layers teal, encoder indigo, decoder plum. Stage 4's two blocks (16 tokens) hold more than any
other component; the decoder's two refine convolutions hold 1,770,240 between them; the spectral
compressor and decompressor are invisible at this scale (0.2\% together).}
\label{fig:params:bars}
\end{figure}

\coursefig[0.95\linewidth]{iw-26-params}{The course's view of the same budget by component, bar widths to
scale: the blocks (51.4\%) and decoder (39.3\%) carry almost all of it.}{fig:params:course}

Five lessons follow from the table and the chart.

\begin{enumerate}
\item \textbf{Stage 4 is the most expensive stage}: 1.6 million parameters, 29\% of the model, while
processing 16 tokens. Parameters follow width squared, not token count.
\item \textbf{The decoder is 39\% of the model.} ``Lightweight'' means light compared with a
segmentation decoder running at full resolution with skip connections, not small.
\item \textbf{Two layers are 32\% of everything.} The refine convolutions hold
$590{,}080 + 1{,}180{,}160 = 1{,}770{,}240$, and $1{,}770{,}240 / 5{,}506{,}629 \approx 0.321$.
\item \textbf{The spectral bottleneck costs 0.2\% and saves far more.} Without it, stage 1's patch
embedding would read 165 channels: $165 \times 32 \times 16 + 32 = 84{,}480 + 32 = 84{,}512$ instead of
16,416, and every activation in the encoder's first convolution would be five times larger.
\item \textbf{Parameters are not compute}, which needs its own section.
\end{enumerate}

\section{Parameters are not compute}\label{sec:params:compute}

\index{multiply-add}\index{compute versus parameters}A layer's work is its multiply-adds per output position times the number of positions. A layer with few
parameters applied at many positions can cost more than a large layer applied at few.

\begin{worked}[label={we:params:macs}]{Multiply-adds of three layers}
For one 128 $\times$ 128 patch, weights per output position (no biases) times positions:
\begin{align*}
\text{compressor:}\quad & 165 \times 32 = 5{,}280 \text{ per pixel}; & 5{,}280 \times 16{,}384 &= 86{,}507{,}520\\
\text{stage-4 fc1:}\quad & 256 \times 1{,}024 = 262{,}144 \text{ per token}; & 262{,}144 \times 16 &= 4{,}194{,}304\\
\text{decoder refine 2:}\quad & 256 \times 512 \times 9 = 1{,}179{,}648; & 1{,}179{,}648 \times 1{,}024 &= 1{,}207{,}959{,}552
\end{align*}
The compressor has 50 times fewer parameters than stage 4's fc1 ($5{,}312$ against $263{,}168$) but does
about 20 times the work, because it runs at 16,384 pixels instead of 16 tokens.
\end{worked}

Hooking every convolution and linear layer of the real model during one unmasked forward pass of a
128 $\times$ 128 patch gives the totals in \cref{tab:params:macs}\tagC{} (attention's own matrix
products, at most a few million per block with 16 keys, are not included).

\begin{table}[htbp]\centering
\caption{Multiply-adds of all convolution and linear layers, one unmasked 128 $\times$ 128 patch,
counted by \texttt{research/textbook/tools/compute\_params.py}.}\label{tab:params:macs}
\begin{tabular}{@{}lrrr@{}}\toprule
Component & Multiply-adds & Share of work & Share of parameters \\\midrule
compressor & 86,507,520 & 3.58\% & 0.10\% \\
stage 1 (embed, blocks, norm) & 42,270,720 & 1.75\% & 3.19\% \\
stage 2 & 29,229,056 & 1.21\% & 4.64\% \\
stage 3 & 44,318,720 & 1.83\% & 16.89\% \\
stage 4 & 31,358,976 & 1.30\% & 35.77\% \\
decoder & 2,096,627,712 & 86.75\% & 39.31\% \\
decompressor & 86,507,520 & 3.58\% & 0.10\% \\\midrule
total & 2,416,820,224 & 100\% & \\\bottomrule
\end{tabular}
\end{table}

\begin{sourcenote}[: where the compute goes]
The course says of \texttt{torchinfo}'s 2.32 G multiply-adds that ``most of that happens at stage 1,
where the parameters are fewest but the tensors are largest''. Counting layer by layer contradicts it:
the whole encoder does about 147 million multiply-adds (6\%), and the \emph{decoder} does 2.10 billion
(87\%), because its fuse and refine convolutions run at all 1,024 positions with 256 to 1,024 channels.
The two refine convolutions alone are 1.81 billion (\cref{ch:decoder}). The lesson the course wanted
(few parameters at many positions can dominate) is right, and the compressor and decompressor show it;
but the layer that dominates is the decoder. The count here is 2.42 billion against \texttt{torchinfo}'s
2.32 (\file{research/model_break_down/05_hyperspectral_segformer_mae.md}, line 293); the 4\% gap is in
how the two tools count, not in the model, and does not change the ranking.
\end{sourcenote}

\section{Parameters a masked pass never trains}\label{sec:params:idle}

\index{idle parameters}There is one more thing a count reveals, and it follows from \cref{ch:encoder}. In a masked pass
stage 1 has fewer than $H \times W$ tokens, so ESA skips its reduction convolution and norm, and Mix-FFN
skips its depthwise convolution. Layers that are skipped get no gradient. Per stage-1 block:
\begin{align*}
\text{reduction conv} + \text{reduction norm} + \text{depthwise} &= 65{,}568 + 64 + 1{,}280 = 66{,}912,\\
\text{two blocks} &= 2 \times 66{,}912 = 133{,}824 .
\end{align*}
That is $133{,}824 / 5{,}506{,}629 \approx 2.4\%$ of the model\tagC. A backward pass through the real
model with a 65\% keep mask confirms it: exactly these twelve tensors, 133,824 numbers, have zero
gradient (\file{research/textbook/tools/compute_params_output.txt}).

\begin{sourcenote}[: idle weights during masked training]
Masked pretraining always passes a keep mask
(\file{app/foundation_models/trainers/hyperspectral_segformer_mae_trainer.py}, line 231), and masked
inference does too, so these 133,824 weights keep their initial values through all of reconstruction
training. They are used, untrained, by any \emph{unmasked} pass: \code{keep_mask=None} in
\code{forward}, and the segmentation head's pass in fine-tuning, which encodes without a keep mask
(\file{spectral_stem_seg_former_mae.py}, \code{segment}), where they first receive a gradient
(\cref{ch:seghead}). The course does not mention this. It is not a bug in the masked path, which never
uses them, but it is a trap for anyone who evaluates a pretrained model unmasked.
\end{sourcenote}

\section{Other models: which figure belongs to which}\label{sec:params:others}

The courses quote several parameter figures. They belong to different models, and mixing them up is
easy.

\subsection{The v0.1.0 model, \texorpdfstring{$D = 24$}{D = 24}}

Shrinking $D$ from 32 to 24 changes four layers: the compressor, stage 1's patch embedding (its input
width), the decompressor, and the last refine convolution (its output is $16D$).
\begin{align*}
\text{compressor:}\quad & 165 \times 24 + 24 + 2 \times 24 = 3{,}960 + 24 + 48 = 4{,}032 && (\text{was } 5{,}376)\\
\text{patch embed 1:}\quad & 24 \times 32 \times 16 + 32 + 64 = 12{,}288 + 96 = 12{,}384 && (\text{was } 16{,}480)\\
\text{decompressor:}\quad & 24 \times 165 + 165 = 3{,}960 + 165 = 4{,}125 && (\text{was } 5{,}445)\\
\text{refine 2:}\quad & 256 \times 384 \times 9 + 384 = 884{,}736 + 384 = 885{,}120 && (\text{was } 1{,}180{,}160)
\end{align*}
The differences: $-1{,}344$, $-4{,}096$, $-1{,}320$, $-295{,}040$, together $-301{,}800$. So
\[
5{,}506{,}629 - 301{,}800 = 5{,}204{,}829
\]
parameters, plus $660 + (24 + 24 + 1) = 709$ buffers, $5{,}205{,}538$ in all: the v0.1.0 manifest's
figure, and the CPU build's. Ninety-eight percent of the saving is the decoder's last convolution:
shrinking $D$ mostly shrinks the decoder.

\subsection{The spectral-stem model (Indradhanu V2, FT04)}

From \cref{ch:crossstem} on, the fixed 165-band compressor is replaced by a wavelength-aware stem, and
fine-tuning adds a segmentation head. \Cref{tab:params:stem} gives that model's split; its encoder and
reconstruction decoder are \emph{the same classes at the same widths} as above, which is why
3,331,072 and 2,164,736 reappear.

\begin{table}[htbp]\centering
\caption{The spectral-stem model with the reduced-attention stem and a 5-class segmentation head, as
configured for FT04\tagM{research/experiments/indradhanu-finetune/FT04-final-joint/config.json}.
Counts from the CPU build\tagM{research/textbook/tools/compute_params_output.txt}; the course quotes
the same split (IW67, IW40, IW47).}\label{tab:params:stem}
\begin{tabular}{@{}llr@{}}\toprule
Component & Where it is taught & Parameters \\\midrule
reduced-attention compressor (\texttt{spectral\_in}) & \cref{ch:specdecoder} & 14,721 \\
BatchNorm after the stem & & 64 \\
encoder & \cref{ch:encoder}, this chapter & 3,331,072 \\
reconstruction decoder & \cref{ch:decoder}, this chapter & 2,164,736 \\
wavelength decoder (\texttt{spectral\_out}) & \cref{ch:specdecoder} & 4,257 \\
segmentation head (width 128, 5 classes) & \cref{ch:seghead} & 371,536 \\\midrule
total & & 5,886,386 \\\bottomrule
\end{tabular}
\end{table}

Two checks. Without the segmentation head the reconstruction model has $5{,}886{,}386 - 371{,}536 =
5{,}514{,}850$ parameters, only 8,221 more than the 165-band model: the new stem and decoder (18,978
together, plus the BatchNorm's 64) replace a compressor and decompressor of 10,821. And the head is the
decoder class of \cref{sec:params:decoder} at width 128 with $5 \times 16 = 80$ output channels:
\begin{align*}
\text{projections:}\quad & (32 + 64 + 160 + 256) \times 128 + 4 \times 128 = 65{,}536 + 512 = 66{,}048\\
\text{fuse:}\quad & 512 \times 128 + 128 = 65{,}536 + 128 = 65{,}664\\
\text{refine 1:}\quad & 128 \times 128 \times 9 + 128 = 147{,}456 + 128 = 147{,}584\\
\text{refine 2:}\quad & 128 \times 80 \times 9 + 80 = 92{,}160 + 80 = 92{,}240\\
\text{head:}\quad & 66{,}048 + 65{,}664 + 147{,}584 + 92{,}240 = 371{,}536 .
\end{align*}
None of these counts contains the band count 202: the stem and the wavelength decoder work for any band
list, which is their point (\cref{ch:ceiling}).

\begin{drift}
As of 2026-10-29 the shipped Indradhanu V2 (v2.1.0) is \emph{Indradhanu Heavy}\index{Indradhanu Heavy!parameter count} (FTH3 epoch 5), a wider
model with a progressive decoder: stage widths 96, 192, 480, 768, decoder width 512, $D = 64$,
segmentation width 256\tagM{research/experiments/indradhanu-finetune/FTH3-heavy-joint-more2/config.json}.
Building it on CPU gives 35,315,314 parameters\tagC: encoder 29,580,288, progressive decoder
4,279,616, segmentation head 1,431,120, stem 17,825 + 128 + 6,337
(\file{research/textbook/tools/compute_params_output.txt}). Without the segmentation head the total is
$35{,}315{,}314 - 1{,}431{,}120 = 33{,}884{,}194$, the figure \cref{ch:heavy} and the pretraining README
quote. \Cref{ch:heavy} tells that story; every
count in this chapter is for the standard widths.
\end{drift}

\begin{confusions}
\confusion{Does the number of heads change the count?}{No. Multi-head attention splits the same four
$C \times C$ matrices into groups; no weights are added. The rules never use it.}
\confusion{Why is the count independent of patch size?}{Convolutions and linear layers reuse the same
weights at every position. A bigger patch means more positions, more work and more memory, not more
parameters.}
\confusion{Are the normalisation statistics parameters?}{No, buffers: saved in the file, never updated
by gradient. They explain the 725 between 5,506,629 and the manifest's 5,507,354.}
\confusion{The decoder is called lightweight; is it small?}{It is 39\% of the parameters and 87\% of the
convolution work. Lightweight compared with full-resolution segmentation decoders, not small.}
\confusion{Is 5.5 million big?}{Small by modern standards (ResNet-50 has about 25 million). It is
deliberately compact: it runs in a worker container, sometimes on a CPU.}
\confusion{Which total is ``Indradhanu's''?}{It depends on the model: 5,506,629 for the 165-band v0.2.0,
5,204,829 for v0.1.0, 5,886,386 for the spectral-stem FT04 model with its head, 35,315,314 for Heavy.
Quote the model with the number.}
\end{confusions}

\begin{checkyourself}
\item\label{ex:params:1} How many parameters does \code{Conv2d(64, 128, k=3)} have? And the depthwise
\code{Conv2d(64, 64, k=3, groups=64)}? How many times fewer?
\item\label{ex:params:2} Use \eqref{eq:params:block} to count one block of width $C = 128$ with $R = 4$,
and check it against a piece-by-piece count.
\item\label{ex:params:3} A colleague proposes three blocks at stage 4 instead of two. How many parameters
does that add, and what fraction of the 165-band model is that?
\item\label{ex:params:4} Compute the 165-band model's total if the decoder width were 128 instead of 256
(everything else unchanged). Which decoder layer shrinks most?
\item\label{ex:params:5} Reconcile 5,506,629 with 5,507,354 from memory of the model's structure, without
looking: which buffers, how many of each?
\item\label{ex:params:6} The compressor has 5,376 parameters and the decoder's fuse convolution 262,400.
For one 128 $\times$ 128 patch, which does more multiply-adds, and by what factor?
\end{checkyourself}

\nextchapter{ch:chakshu}
